% ===== PRÉAMBULE CANONIQUE — à coller VERBATIM en tête de CHAQUE .tex =====
\documentclass[11pt,a4paper]{article}
\usepackage[utf8]{inputenc}\usepackage[T1]{fontenc}\usepackage{lmodern}
\usepackage[french]{babel}
\usepackage{amsmath,amssymb,mathtools}\usepackage{icomma}
\usepackage{siunitx}\sisetup{locale=FR}  % \qty{}{}, \si{}, \ang{}
\usepackage{xcolor}\usepackage{colortbl}
\usepackage{tikz}\usepackage{pgfplots}\pgfplotsset{compat=1.18}
\usepackage[siunitx,europeanresistors]{circuitikz} % circuits (élec.)
\usepackage[most]{tcolorbox}\usepackage{enumitem}\usepackage{array,booktabs}
\usepackage[a4paper,margin=2.2cm]{geometry}
\usetikzlibrary{arrows.meta,calc,angles,quotes,positioning,patterns,%
  backgrounds,shapes.geometric,decorations.pathmorphing,decorations.markings,intersections}
\definecolor{primary}{RGB}{222,49,99}\definecolor{primarydark}{RGB}{150,22,55}
\definecolor{defcol}{RGB}{0,114,178}\definecolor{methcol}{RGB}{52,92,148}
\definecolor{excol}{RGB}{0,135,90}\definecolor{attncol}{RGB}{200,40,55}
\tcbset{boxbase/.style={enhanced,breakable,boxrule=0.9pt,arc=2.5pt,
  left=3mm,right=3mm,top=2.3mm,bottom=2.3mm,fonttitle=\bfseries,coltitle=white}}
% --- boîtes TUTORIEL (noms EXACTS, ne pas renommer) ---
\newtcolorbox{cle}[1][]{boxbase,colback=defcol!6,colframe=defcol,
  title={\textsc{À retenir}\ifx\relax#1\relax\relax\else\textmd{ : \itshape #1}\fi}}
\newtcolorbox{methode}[1][]{boxbase,colback=methcol!7,colframe=methcol,sharp corners,
  title={\textsc{Méthode}\ifx\relax#1\relax\relax\else\textmd{ --- \itshape #1}\fi}}
\newtcolorbox{exemplebox}[1][]{boxbase,colback=excol!5,colframe=excol,
  title={\textsc{Exemple}\ifx\relax#1\relax\relax\else\textmd{ : \itshape #1}\fi}}
\newtcolorbox{piege}[1][]{boxbase,colback=attncol!6,colframe=attncol,
  title={\textsc{Piège}\ifx\relax#1\relax\relax\else\textmd{ : \itshape #1}\fi}}
\newtcolorbox{remarque}[1][]{boxbase,colback=black!4,colframe=black!45,
  title={\textsc{Remarque}\ifx\relax#1\relax\relax\else\textmd{ : \itshape #1}\fi}}
% --- boîtes EXERCICE / CORRIGÉ (noms EXACTS, auto-comptées) ---
\newtcolorbox[auto counter]{exo}[1][]{boxbase,colback=primary!4,colframe=primary,
  title={\textsc{Exercice}~\thetcbcounter\ifx\relax#1\relax\relax\else\textmd{ --- \itshape #1}\fi}}
\newtcolorbox[auto counter]{corrige}[1][]{boxbase,colback=excol!4,colframe=excol,
  title={\textsc{Corrigé}~\thetcbcounter\ifx\relax#1\relax\relax\else\textmd{ --- \itshape #1}\fi}}
\newcommand{\vv}[1]{\vec{#1}}\newcommand{\ds}{\displaystyle}
\newcommand{\R}{\mathbb{R}}\newcommand{\N}{\mathbb{N}}\newcommand{\Z}{\mathbb{Z}}
% =========================================================================
\begin{document}
\sisetup{per-mode=symbol}

\begin{corrige}[le parachutiste : vrai ou faux ?]
On raisonne avec $a=g-\dfrac{F_f}{m}$ (sens de la chute).
\begin{enumerate}[label=\Alph*.]
  \item \textbf{Faux}. Si $F_f>G$, la résultante est dirigée vers le haut : le parachutiste \emph{ralentit}, sa vitesse diminue.
  \item \textbf{Faux}. Si $F_f<G$, la résultante est vers le bas : il \emph{accélère}, la vitesse n'est pas constante.
  \item \textbf{Vrai}. Si $F_f=G$, la résultante est nulle : $a=0$, vitesse constante (vitesse limite).
  \item \textbf{Vrai}. Avec de l'air, $F_f>0$ donc $a=g-\tfrac{F_f}{m}<g$ ; $g$ est le maximum (chute libre).
\end{enumerate}
\end{corrige}

\begin{corrige}[trois ballons qui tombent]
L'affirmation \textbf{incorrecte est C}.
\begin{itemize}[leftmargin=1.4em,topsep=2pt,itemsep=1.5pt]
  \item \textbf{C (fausse)} : l'accélération vaut $a=g-\tfrac{F_f}{m}$, toujours $<g$. Elle ne peut donc
        \emph{jamais} dépasser $g$.
  \item Les autres sont vraies : A (quand $v\nearrow$, $F_f\nearrow$ donc $a\searrow$) ; B (si $F_f=G$ alors $a=0$) ;
        D (le frottement croît avec la vitesse) ; E (à même instant, $\tfrac{F_f}{m}$ est plus petit pour la
        grande masse, donc $a$ plus grand : $a_C>a_B>a_A$).
\end{itemize}
\end{corrige}

\begin{corrige}[les trois phases d'un saut]
Toujours $a=g-\dfrac{F_f}{m}$.
\begin{enumerate}[label=\alph*)]
  \item $F_f\approx 0 \Rightarrow a\approx g=\qty{10}{\metre\per\second\squared}$.
  \item $a=10-\dfrac{600}{80}=10-\num{7,5}=\qty{2,5}{\metre\per\second\squared}$.
  \item Vitesse constante $\Rightarrow a=0 \Rightarrow F_f=G=mg=80\times 10=\qty{800}{\newton}$.
\end{enumerate}
\end{corrige}

\begin{corrige}[masse, poids et chute : le vrai du faux]
Correctes : \textbf{2 et 3}.
\begin{enumerate}[label=\arabic*.]
  \item \textbf{Faux}. La masse est invariable ; elle ne dépend pas du lieu.
  \item \textbf{Vrai}. $G=mg$ varie avec $g$ (lieu), la masse non.
  \item \textbf{Vrai}. En chute libre $a=g$ pour tous : ils touchent le sol ensemble.
  \item \textbf{Faux}. L'accélération de chute libre ne dépend pas de la masse ($a=g$).
  \item \textbf{Faux}. Avec frottement $a=g-\tfrac{F_f}{m}<g$ : jamais au-dessus de $g$.
\end{enumerate}
\end{corrige}

\begin{corrige}[la corde en chute libre]
\begin{enumerate}[label=\alph*)]
  \item Sans air, chaque corps n'est soumis (essentiellement) qu'à son poids : $a=g$ pour A comme pour B.
  \item La force sur la corde est \textbf{nulle}. Les deux corps « tombent » avec la même accélération $g$ :
        B n'a pas tendance à s'éloigner de A ni à le tirer, la corde ne se tend pas. (Si la corde était
        coupée, rien ne changerait à leur mouvement.)
\end{enumerate}
\end{corrige}
\end{document}
